\chapter*{Abstract}
\addcontentsline{toc}{chapter}{Abstract} % Añade "Abstract" a la Tabla de Contenido

Producing the pulse an application requires from an ultrafast laser amounts to fixing the
parameter configuration of its cavity, a black-box, non-convex and multistable optimization
problem for which traditional tuning methods fall short. The literature has responded with
nature-inspired algorithms and machine-learning techniques, but leaves two gaps: it has
concentrated almost entirely on the passive mode-locking regime, leaving the active one
underexplored, and it is organized as a collection of isolated case studies, without
comparisons that hold the laser, the parameter space and the objective function fixed while
varying only the algorithm. This thesis addresses both gaps through a controlled comparative
study of three gradient-free strategies —a genetic algorithm (GA), a hybrid with a neural
surrogate (NN+GA), and particle swarm optimization (PSO)— on a Haus master-equation simulator
validated against analytical solutions, across the active and passive regimes, under a single
merit function and common metrics of quality, cost and robustness.

Over a matrix of sixty independent runs with non-parametric statistical testing, the three
strategies reach feasible solutions, but the preferable method depends on the regime. In the active regime, PSO significantly dominates in quality, dispersion,
convergence speed and optimum robustness, associated with stronger soliton shortening; in the passive regime, PSO and GA are statistically
tied. The neural surrogate matches the GA in the active regime using $5{.}7$ times fewer
simulator evaluations and sacrifices quality in the passive one, although a paired control
shows that, at that budget, its contribution cannot be distinguished from that of its own
initial sample. A single factor —the shape of each regime's optimization landscape,
regular in the active case and multistable in the passive one— accounts for all three
findings, and local robustness emerges as a comparison dimension that discriminates between
methods when quality alone does not. The work, carried out in simulation, contributes a
reproducible framework and evidence that the active regime is, for this family of methods, the
more tractable one; experimental validation and the extension to deep reinforcement learning
are left as future work.

\vspace{1em}
\noindent\textbf{Keywords:} ultrafast lasers, mode-locking, gradient-free optimization,
genetic algorithm, particle swarm optimization, neural surrogate models, Haus master
equation, machine learning.
